% ══════════════════════════════════════════════════════════════════════
% Viewpoint 4 Corrigendum: Cercis-Fisher is Conditional Equivalence, NOT Natural Emergence
% 4Cercis-Fisher
%
% This is a standalone corrigendum to be inserted after Section 4 (Information-Geometric
% Bulk-Boundary Correspondence) of gauge_formalized.tex, or read alongside it.
%
% Based on formal peer review (docs/reviews/GAUGE_5VIEWPOINTS_FINAL.md,
% docs/reviews/GAUGE_VIEWPOINTS_REVIEW.md).
% ══════════════════════════════════════════════════════════════════════

\documentclass{article}
\usepackage[UTF8]{ctex}
\usepackage{amsmath,amssymb,amsthm}
\usepackage{hyperref}
\usepackage{enumitem}
\usepackage[margin=2.5cm]{geometry}

% ─── Math notation (consistent with gauge_formalized.tex) ───
\newcommand{\GeoDist}{\mathrm{GeoDist}}
\newcommand{\KL}{\mathrm{KL}}
\newcommand{\Fisher}{\mathcal{I}}
\newcommand{\cercis}{\mathcal{C}}
\newcommand{\cercisnorm}{\bar{\cercis}}
\newcommand{\Situs}{\mathcal{S}}
\newcommand{\ExpFam}{\mathcal{E}}
\newcommand{\Od}{\mathrm{O}(d)}
\newcommand{\SOthree}{\mathrm{SO}(3)}
\newcommand{\ZBetti}{\beta_1}
\newcommand{\Lone}{L_1}
\newcommand{\edgs}{\mathcal{E}}
\newcommand{\verts}{\mathcal{V}}
\newcommand{\face}{\mathcal{F}}
\newcommand{\grph}{\mathcal{G}}
\newcommand{\Gflatquot}{\mathcal{M}_{\mathrm{flat}}}
\newcommand{\triv}{\mathbf{1}}
\newcommand{\doneT}{d_1^{\mathsf{T}}}
\DeclareMathOperator{\im}{im}
\DeclareMathOperator{\Log}{Log}
\DeclareMathOperator{\Ad}{Ad}
\DeclareMathOperator{\Hom}{Hom}
\DeclareMathOperator{\Stab}{Stab}
\newcommand{\E}{\mathbb{E}}
\DeclareMathOperator{\dist}{dist}

\title{\textbf{4Corrigendum} \\
\large Cercis-Fisher \\
\large Cercis-Fisher Equivalence: Conditional Equivalence, Not Natural Emergence}
\author{SCX  / SCX Theory Architecture Group}
\date{\today}

\begin{document}
\maketitle

% ══════════════════════════════════════════════════════════════════════
\section{ / Review Verdict}
% ══════════════════════════════════════════════════════════════════════

\begin{center}
\begin{tabular}{|p{6cm}|p{8cm}|}
\hline
\textbf{ / Original Claim} & \textbf{ / Verdict} \\
\hline
Cercis${}^2 \propto$ KL $=$ Fisher

& ⚠️ ``'' / Mathematically correct but ``natural emergence'' is overstated \\
\hline
\end{tabular}
\end{center}

\subsection{ / Specific Issues}

\begin{enumerate}[label=(\arabic*)]
    \item \textbf{``KL''``KL''}
    CercisKL
    $\cercis \approx 2 \cdot \min_{h} \KL(P_A \| P_{d_0 h})$
    ``KL''

    \textbf{Not ``average KL'' but ``minimum KL''.}
    Cercis is the \emph{minimum} KL divergence to the pure-gauge submanifold,
    not an average.

    \item \textbf{``''}
    $\GeoDist^2 \leftrightarrow 2\KL$ 8\ref{sec:assumptions}
    SCX

    \textbf{Not ``natural emergence'' but conditional equivalence.}
    The $\GeoDist^2 \leftrightarrow 2\KL$ equivalence depends on 8 strong assumptions (see \S\ref{sec:assumptions}).
    If any single assumption fails, the equivalence breaks. SCX data is extremely unlikely to form an exponential family.

    \item \textbf{}
    $\GeoDist^2$$2\KL$Amari-Chentsov
    SCX

    \textbf{Complete breakdown outside exponential families.}
    For general statistical manifolds, the difference between $\GeoDist^2$ and $2\KL$ is controlled
    by the Amari-Chentsov third-order tensor. This tensor is generically non-zero for SCX data,
    so the equivalence degrades.
\end{enumerate}

% ══════════════════════════════════════════════════════════════════════
\section{ / The Eight Assumptions}
\label{sec:assumptions}
% ══════════════════════════════════════════════════════════════════════

 $\GeoDist(p,q)^2 = 2\KL(p\|q) + O(\|p-q\|^3)$ 


For $\GeoDist(p,q)^2 = 2\KL(p\|q) + O(\|p-q\|^3)$ to hold rigorously,
the following eight assumptions must be simultaneously satisfied.
Failure of any single one degrades the equality.

\vspace{1em}

\begin{enumerate}[label=\textbf{H\arabic*:}, leftmargin=*, itemsep=0.8em]

    \item \textbf{ (Exponential Family Assumption).} \\
    \textbf{}  $\mathcal{P} = \{p(x \mid \theta)\}$ 
     $T(x)$  $\psi(\theta)$ 
    $p(x \mid \theta) = h(x) \exp(\langle\theta, T(x)\rangle - \psi(\theta))$
    Fisher $\psi$ Hessian
    KLBregman4.3 / Lemma~\ref{lem:kl_bregman}
    \\[0.3em]
    \textbf{English:} The output distribution family $\mathcal{P}$ must be an exponential family.
    Only then is the Fisher matrix the Hessian of $\psi$, and KL divergence equals Bregman divergence.
    \\[0.3em]
    \textbf{} $\KL(\theta_1\|\theta_2)$  $\psi$ Bregman
    Taylor $g_{ij}\Delta\theta^i\Delta\theta^j$
    $\GeoDist^2$  $2\KL$  Amari-Chentsov 6.2 / Open Problem~\ref{prob:non_exp_fam}
    \\[0.3em]
    \textbf{Consequence of failure:} Outside exponential families, KL is no longer a Bregman divergence,
    and the second-order Taylor term is no longer $g_{ij}\Delta\theta^i\Delta\theta^j$.
    The difference is controlled by the Amari-Chentsov third-order tensor (Open Problem~\ref{prob:non_exp_fam}).

    \item \textbf{ (Small-Distance / Local Approximation Assumption).} \\
    \textbf{} $\theta_p$  $\theta_q$ 
    $O(\|\theta_q - \theta_p\|^3)$ Taylor
    \\[0.3em]
    \textbf{English:} $\theta_p$ and $\theta_q$ must be sufficiently close so that third-order
    and higher remainders $O(\|\theta_q - \theta_p\|^3)$ are negligible. All derivations
    truncate the Taylor expansion at second order.
    \\[0.3em]
    \textbf{} large discrepancy
    $\GeoDist^2 \approx 2\KL$ SCX
    \\[0.3em]
    \textbf{Consequence of failure:} For large discrepancies (common in cross-domain SCX experts),
    third-order error terms dominate and $\GeoDist^2 \approx 2\KL$ no longer holds.

    \item \textbf{$C^2$  ($C^2$ Differentiability Assumption).} \\
    \textbf{}  $\theta \mapsto p(x \mid \theta)$  $x$  $C^2$ 
    TaylorFisher
    \\[0.3em]
    \textbf{English:} The map $\theta \mapsto p(x \mid \theta)$ must be $C^2$ for almost all $x$.
    This is required for second-order Taylor expansion and for the Fisher matrix as a second derivative.
    \\[0.3em]
    \textbf{} ReLUFisher
    
    \\[0.3em]
    \textbf{Consequence of failure:} If the density has non-differentiable points (e.g., piecewise
    distributions from ReLU activations), the Fisher metric is undefined at those points and
    geodesics are not uniquely determined.

    \item \textbf{- (Integration-Differentiation Interchangeability Assumption).} \\
    \textbf{} $\partial_\theta \int p(x|\theta) d\mu(x) = \int \partial_\theta p(x|\theta) d\mu(x)$
    Fisher $\E[\partial_\theta \log p] = 0$
    \\[0.3em]
    \textbf{English:} The interchange $\partial_\theta \int p = \int \partial_\theta p$ must hold.
    This is the standard regularity condition ensuring $\E[\partial_\theta \log p] = 0$, which
    underpins the Fisher information definition.
    \\[0.3em]
    \textbf{} Fisher
    \\[0.3em]
    \textbf{Consequence of failure:} If the support depends on $\theta$ (e.g., uniform distributions),
    the standard Fisher information form breaks down.

    \item \textbf{Fisher (Fisher Full-Rank Assumption).} \\
    \textbf{} Fisher $\Fisher(\theta)$ 
     $\Theta$ Fisher——
    \\[0.3em]
    \textbf{English:} The Fisher information matrix $\Fisher(\theta)$ must be positive-definite
    (full rank) at every point. This is required for $(\Theta, g)$ to be a genuine Riemannian
    manifold — a degenerate metric does not define unique geodesics.
    \\[0.3em]
    \textbf{} Fisher
    $\GeoDist$ 
    \\[0.3em]
    \textbf{Consequence of failure:} If parameters are functionally dependent (e.g., normalization
    constraints causing effective dimensionality reduction), the Fisher matrix is singular,
    geodesics may not exist or be unique, and $\GeoDist$ is undefined.

    \item \textbf{ (Small-Connection Approximation Assumption).} \\
    \textbf{}  $a_e$  $h_v$ 
    $\max_e \dist_{\Od}(A_e, I_d) < r_0$ $r_0$ 
     $\mathfrak{so}(d)$ 
    \\[0.3em]
    \textbf{English:} Edge assignments $a_e$ and gauge transformations $h_v$ must be within the
    small-connection limit ($\max_e \dist_{\Od}(A_e, I_d) < r_0$, the matrix logarithm convergence radius).
    Only then does non-abelian gauge-fixing reduce to linear Hodge theory on $\mathfrak{so}(d)$.
    \\[0.3em]
    \textbf{} $90^\circ$Cartan
    $\GeoDist$ 
    \\[0.3em]
    \textbf{Consequence of failure:} Under large rotations (e.g., $90^\circ$ frame misalignment
    between cross-modal experts), nonlinear Cartan terms cannot be neglected, and the relationship
    between $\GeoDist$ and Euclidean distance breaks.

    \item \textbf{ (Euclidean Approximation / Degenerate Fisher Assumption).} \\
    \textbf{} Cercis$\|r\|^2 = \|r_{\text{harm}}\|^2 + \|r_{\text{coexact}}\|^2$
    FisherFisher$\Fisher_{ij}(\theta) = \delta_{ij}$
    
    \\[0.3em]
    \textbf{English:} For the Euclidean Cercis definition ($\|r\|^2 = \|r_{\text{harm}}\|^2 + \|r_{\text{coexact}}\|^2$)
    to coincide with the Fisher definition, the Fisher metric must degenerate to Euclidean:
    $\Fisher_{ij}(\theta) = \delta_{ij}$. This means all directions on the information manifold
    have equal weight and zero curvature.
    \\[0.3em]
    \textbf{} SCXFisher
    CercisFisher
    \\[0.3em]
    \textbf{Consequence of failure:} Real SCX parameter spaces almost certainly have non-trivial
    Fisher curvature. A non-negligible curvature correction exists between Euclidean Cercis and
    Fisher geodesic distance.

    \item \textbf{ (Distribution Compactness / Bounded Support Assumption).} \\
    \textbf{}  $\Theta$ Taylor
    
    \\[0.3em]
    \textbf{English:} All expert output distributions must be sufficiently concentrated on a compact
    subset of $\Theta$, so that the local Taylor expansion is uniformly valid across all
    boundary observations and gauge bulk states.
    \\[0.3em]
    \textbf{} -t
    
    \\[0.3em]
    \textbf{Consequence of failure:} Heavy-tailed distributions (e.g., Student's-t) may place
    some observations outside the local validity domain, causing uncontrolled growth of
    third-order error terms on individual edges.
\end{enumerate}

% ══════════════════════════════════════════════════════════════════════
\section{ / Corrected Claim}
% ══════════════════════════════════════════════════════════════════════

\subsection{ (Chinese Correction)}

{\bf }
\begin{quote}
``CercisFisherCercis${}^2 \propto$ KL
$=$ Fisher——''
\end{quote}

{\bf }
\begin{quote}
{\bf }Cercis{\bf }KL
Fisher
{\bf }——``''{\bf }
Situs
{\bf Amari-Chentsov6.2}
\end{quote}



\vspace{0.5em}
\fbox{%
\begin{minipage}{0.94\textwidth}
\vspace{0.8em}
\noindent {\bf Cercis}
\label{thm:cercis_conditional_equivalence}

{\bf } 8(H1--H8)

{\bf }
\[
\cercis = \min_{d_0 h \in \im(d_0)} \GeoDist(P_A,\; P_{d_0 h})^2
       = 2 \min_{d_0 h \in \im(d_0)} \KL(P_A \| P_{d_0 h}) + O(\|\mathfrak{a}\|^3).
\]

Cercis``''
Fisher

{\bf }
\begin{enumerate}[label=(\roman*)]
    \item (H1)$\GeoDist^2$$2\KL$
    Amari-Chentsov$T_{ijk} = \E[\partial_i \log p \cdot \partial_j \log p \cdot \partial_k \log p]$
    $T_{ijk}$SCX\ref{prob:non_exp_fam}
    \item (H2)
    \item ``''——
    
\end{enumerate}
\vspace{0.5em}
\end{minipage}%
}
\vspace{0.5em}

\vspace{1em}

\subsection{English Correction}

{\bf Original claim (overstated, to be removed):}
\begin{quote}
``The Cercis Score is the geodesic distance under the Fisher information metric.
Under the exponential family assumption, Cercis${}^2 \propto$ mean KL divergence
$=$ squared Fisher geodesic length. The audit score is not artificially defined —
it naturally emerges from the information geometry of the data manifold.''
\end{quote}

{\bf Corrected claim (conditional equivalence):}
\begin{quote}
{\bf Under the exponential family assumption and local approximation,} the Cercis Score
equals twice the {\bf minimum} KL divergence from the observed edge distribution to the
pure-gauge submanifold, which in turn equals (to second order) the squared Fisher geodesic
distance. This provides an information-geometric {\bf interpretation} — not a ``natural emergence''
but a {\bf conditional equivalence} that holds when and only when the Situs manifold
admits an exponential family structure.
{\bf For non-exponential-family data, the relationship degrades by the Amari-Chentsov
tensor (Open Problem~6.2).}
\end{quote}

Explicit formulation:

\vspace{0.5em}
\fbox{%
\begin{minipage}{0.94\textwidth}
\vspace{0.8em}
\noindent {\bf Theorem (Information-Geometric Conditional Equivalence of Cercis, Revised)}
\label{thm:cercis_conditional_equivalence_en}

{\bf Assumptions:} All 8 assumptions (H1--H8) hold simultaneously.

{\bf Conclusion:}
\[
\cercis = \min_{d_0 h \in \im(d_0)} \GeoDist(P_A,\; P_{d_0 h})^2
       = 2 \min_{d_0 h \in \im(d_0)} \KL(P_A \| P_{d_0 h}) + O(\|\mathfrak{a}\|^3).
\]

Cercis measures the minimal Fisher information distance from the observed expert edge
distributions to the ``pure-gauge bulk states'' — the hypothetical world where all
experts can be perfectly aligned.

{\bf Limitations:}
\begin{enumerate}[label=(\roman*)]
    \item The equivalence is rigorous within exponential families (H1); outside them,
    the difference between $\GeoDist^2$ and $2\KL$ is controlled by the Amari-Chentsov
    third-order tensor $T_{ijk} = \E[\partial_i \log p \cdot \partial_j \log p \cdot \partial_k \log p]$,
    which is generically non-zero for SCX data (Open Problem~\ref{prob:non_exp_fam}).
    \item The equivalence holds under local approximation (H2); for large discrepancies,
    third-order terms dominate.
    \item ``Natural emergence'' is an incorrect characterization — this is a
    \emph{constructed equivalence} under strong assumptions, not an intrinsic
    geometric inevitability of the data.
\end{enumerate}
\vspace{0.5em}
\end{minipage}%
}
\vspace{0.5em}

% ══════════════════════════════════════════════════════════════════════
\section{ / Five Breakdown Points}
% ══════════════════════════════════════════════════════════════════════

SCX

Below, we summarize five principal ways the equivalence can break on real SCX data:

\vspace{0.5em}

\begin{enumerate}[label=\textbf{B\arabic*:}, leftmargin=*, itemsep=0.8em]

    \item \textbf{ / Non-Exponential-Family Failure (Most Severe).} \\
    SCXMoEsoftmax
    $T(x)$
    6.2\ref{prob:non_exp_fam}
    \\[0.3em]
    SCX expert outputs (MoE routing weights, softmax scores, attention distributions)
    are extremely unlikely to form an exponential family. Exponential families require specific
    algebraic structure (linear sufficient statistics $T(x)$) that deep learning outputs do not
    generally satisfy. The document itself lists this as Open Problem~6.2 (\ref{prob:non_exp_fam}).

    \item \textbf{``KL''``KL'' / Confusing ``min KL'' with ``average KL''.} \\
    Cercis $\min_h \KL(P_A \| P_{d_0 h})$ KL
    
    \\[0.3em]
    Cercis solves $\min_h \KL(P_A \| P_{d_0 h})$ — a minimization problem, not an average of
    KL divergences over all expert pairs. These are mathematically distinct.

    \item \textbf{ / Hidden Euclidean Approximation Requirement.} \\
    CercisFisher $\Fisher_{ij} = \delta_{ij}$ (H7)
    ——SCX
    \\[0.3em]
    The equivalence between Euclidean-norm Cercis and Fisher geodesic distance requires
    $\Fisher_{ij} = \delta_{ij}$ (H7) — meaning the information manifold must be everywhere flat,
    which is extremely unlikely on real SCX parameter spaces.

    \item \textbf{4.2 / Intrinsic Correction in Theorem~4.2 Proof.} \\
    $\im(d_0)$ 
     $\ker(\Lone)$\ref{thm:cercis_bulk_boundary}
    \\[0.3em]
    Choosing the wrong submanifold loses the harmonic component. The correct definition is:
    the bulk submanifold is parameterized by exact connections (pure gauge) $\im(d_0)$,
    not harmonic bulk states $\ker(\Lone)$. Theorem~\ref{thm:cercis_bulk_boundary} identifies
    and corrects this issue in its proof.

    \item \textbf{ / Locality Constraint.} \\
     $\|\theta_q - \theta_p\|$ 
     $O(\|\Delta\theta\|^3)$ 
    \\[0.3em]
    All results hold only when $\|\theta_q - \theta_p\|$ is sufficiently small.
    When inter-configuration expert differences are large (e.g., different domains, different
    training stages), the third-order error term $O(\|\Delta\theta\|^3)$ is non-negligible.

\end{enumerate}

% ══════════════════════════════════════════════════════════════════════
\section{ / Recommended Revisions}
% ══════════════════════════════════════════════════════════════════════

\subsection{ / Changes Needed in gauge\_formalized.tex}

\begin{enumerate}[label=(\arabic*)]

    \item \textbf{234--238}
    ``''
    \begin{quote}
        ``FisherKL...''
    \end{quote}
    
    \begin{quote}
        ``FisherKL...''
    \end{quote}
    

    \item \textbf{\ref{sec:info}}
    
    \begin{quote}
        ``FisherKL / Fisher Geodesic--KL Equivalence under Exponential Families''
    \end{quote}
    
    \begin{quote}
        ``FisherKL / Conditional Fisher Geodesic--KL Equivalence under Exponential Families''
    \end{quote}

    \item \textbf{\ref{thm:cercis_bulk_boundary}}
    H1--H8

    \item \textbf{\ref{thm:unified}}
    
    \begin{quote}
        ``CercisFisher$2\KL$''
    \end{quote}
    
    \begin{quote}
        ``CercisFisher
        $2\KL$6.2''
    \end{quote}

    \item \textbf{1529--1532}
    
    \begin{quote}
        ``Cercis: $\min_{d_0 h} \GeoDist(P_A, P_{d_0 h})^2$'' 
        ``Cercis: $\approx 2\min_{d_0 h} \KL(P_A \| P_{d_0 h})$''
    \end{quote}
    ``''

    \item \textbf{\ref{sec:comparison}``''1560--1564}
    
    \begin{quote}
        ``SCX$\Situs$...''
    \end{quote}
    
    \begin{quote}
        ``SCX$\Situs$
        6.2Amari-Chentsov''
    \end{quote}

    \item \textbf{31678--1684}
    
    \begin{quote}
        ``FisherKL...Cercis''
    \end{quote}
    
    \begin{quote}
        ``FisherKL...Cercis8''
    \end{quote}

\end{enumerate}

% ══════════════════════════════════════════════════════════════════════
\section{6.2 / Cross-Reference to Open Problem 6.2}
% ══════════════════════════════════════════════════════════════════════

\ref{prob:non_exp_fam}``Fisher-KL''


\begin{quote}
SCXFisherKL
Amari-Chentsov $T_{ijk} = \E[\partial_i \log p \cdot \partial_j \log p \cdot \partial_k \log p]$
SCX
\end{quote}

This corrigendum is directly linked to Open Problem~\ref{prob:non_exp_fam}
(``Fisher-KL Equivalence Outside Exponential Families''), which states:

\begin{quote}
For general SCX output distributions (potentially non-exponential-family),
the relationship between Fisher geodesic distance and KL divergence is controlled
by the Amari-Chentsov third-order tensor $T_{ijk} = \E[\partial_i \log p \cdot \partial_j \log p \cdot \partial_k \log p]$.
Bounds on this tensor for empirical SCX output distributions are needed.
\end{quote}

\vspace{1em}
\noindent\rule{\textwidth}{0.5pt}
\vspace{0.5em}

\noindent {\small \textbf{} 
\texttt{docs/reviews/GAUGE\_5VIEWPOINTS\_FINAL.md}``''
\texttt{docs/reviews/GAUGE\_VIEWPOINTS\_REVIEW.md}
4$\checkmark$``''``''}

\noindent {\small \textbf{Review basis:} This corrigendum is based on:
\texttt{docs/reviews/GAUGE\_5VIEWPOINTS\_FINAL.md} (verdict: ``overstated'') and
\texttt{docs/reviews/GAUGE\_VIEWPOINTS\_REVIEW.md} (detailed analysis).
The verdict on Viewpoint~4 is: $\checkmark$ mathematically correct relationship but
``natural emergence'' is severely overstated — downgrade to ``conditional equivalence''.}

\end{document}
